\section{Related Work}

\textbf{Probability calibration.} Producing well-calibrated
probabilities from classifier scores is a long-studied
problem~\cite{niculescu2005predicting}, with post-hoc maps including
Platt scaling~\cite{platt1999probabilistic} and isotonic
regression~\cite{zadrozny2002transforming}, and renewed attention to
the miscalibration of high-capacity neural
models~\cite{guo2017calibration}. The present paper does not propose a
new calibration map; it studies the temporal-update layer that sits
above any such map, asking how the calibration target should be
aggregated across past windows. The finding, that a stale or
over-smoothed target undoes the benefit of recalibration, is
orthogonal to and compatible with the choice of map.

\textbf{Concept drift.} The failure mode we document is a concept-drift
phenomenon: the distribution shifts between the window the
recalibrator is fit on and the window it is deployed on. The drift
literature surveyed by Gama et al.~\cite{gama2014survey} catalogues
exactly the adaptive-window and forgetting-factor mechanisms that our
results, by elimination, point toward as the right class of fix.

\textbf{Vulnerability prioritization and exploit likelihood.}
Operational prioritization spans severity scoring with
CVSS~\cite{mell2007cvss}, stakeholder-specific categorization with
SSVC~\cite{spring2021ssvc}, empirical exploitation studies relating
severity to observed exploits~\cite{allodi2014comparing}, and
data-driven exploit-likelihood prediction with
EPSS~\cite{jacobs2021epss}. Our scorer consumes EPSS as one input and
inherits its non-stationarity; the calibration-update question we
study applies to any of these scores when it is recalibrated online
against an evolving fleet.

\textbf{Lag-1 online calibration baseline.}
The online-calibration study~\cite{paper7online} introduced the rolling-history grid-search
recalibration procedure and reported the capacity-conditional hazard
that motivates this paper. The present sweep inherits the online-calibration study's
\textsf{fixed}/\textsf{lag1}/\textsf{offline} strategies and adds the
two smoothing variants.

\textbf{High-capacity collapse.} our earlier capacity-decay study
identified the high-capacity regime where HygienePrio-full collapses
even with fixed weights. The present paper does not attempt to rescue
the absolute floor (that question is taken up in the future-work
section); instead it asks whether a smoothed online recalibration
procedure can match or exceed the fixed-weight baseline at the
collapse cells.

\textbf{Exponential smoothing.}
The EWMA estimator dates back at least to Holt's seasonal-forecasting
work~\cite{holt2004ewma}; in machine-learning settings it appears
under the name ``geometric averaging'' as a low-variance bias-bearing
substitute for sliding-window means. The use here, weighting a
finite past history of grid-search calibration scores, is a direct
adaptation; no novel estimator is claimed.

\textbf{Multi-paper sequence.}
A connected line of our prior work, spanning the HygieneBench synthetic
benchmark, the HygienePrio scorer~\cite{paper4hygieneprio}, a temporal-stability
study across rolling maintenance windows, a capacity-indexed decay analysis,
and the online-calibration study~\cite{paper7online}, together defines the
methodological substrate within which this paper sits. This paper
extends the calibration-strategy dimension by one step (single-lag $\to$
finite-history weighted aggregation).

\textbf{Policy mandates.} CISA BOD~22-01~\cite{cisa2021bod2201},
23-01~\cite{cisa2022bod2301}, NIST SP~800-40 Rev.~4~\cite{nist2022sp80040r4},
and the 2026 DBIR~\cite{verizon2026dbir} establish the operational
scheduling regime under which any deployable recalibration procedure
must operate.
